\section{Failure Analysis}

\subsection{Why analytic success does not imply raw-sample success}

The analytic spectrum-space experiment deliberately removes the hardest part
of the problem: estimating \(\zeta(q)\) from a finite sample. Its success shows
that the spectral geometry is learnable. It does not show that the same
geometry will be visible in a short raw time series. The failure-attribution
experiment closes this gap by comparing analytic target spectra,
deterministic estimated spectra, and neural empirical spectra for the same
process families. The observed failure occurs primarily between the raw sample
and the empirical spectrum.

This distinction is important for interpretation. If the analytic calibrator
failed, the diagnostic design would be questionable. If deterministic
estimators succeeded but the neural encoder failed, the issue would be neural
training. Instead, deterministic and neural empirical spectra both struggle to
preserve MRW curvature under the tested short-window conditions. The evidence
therefore points to finite-sample spectral identifiability as the bottleneck.

\subsection{Fake curvature and true curvature are coupled}

Raw zeta alignment shows a useful but incomplete success: it suppresses fake
curvature in fGn and Gaussian controls. This reduces false MRW compatibility.
However, the same smoothing and alignment pressure can suppress true MRW
curvature. Curvature-preserving alignment was introduced to counter this
effect, using band-specific MRW curvature losses and third-difference
smoothness. Even then, medium/high-\(\lambda^2\) MRW curvature was not
reliably restored.

This trade-off is not just an implementation detail. In short samples, small
quadratic curvature and finite-sample noise can be comparable in magnitude.
Suppressing noise risks flattening real curvature; preserving curvature risks
letting monofractal noise appear multifractal. This is precisely why the final
framework reports boundary and instability diagnostics instead of forcing a
hard MRW/non-MRW decision.

\subsection{What the negative result means}

The negative identifiability result should not be read as a failure of
multifractal analysis in general. It is a statement about the tested finite
regime: short windows, limited scale support, the chosen \(q\)-grid, and the
implemented structure-function estimator family. Larger samples, richer scale
ranges, or alternative estimators such as wavelet leaders may improve recovery
in some settings. The present result says that, under the conditions studied
here, \(\lambda^2_{\mathrm{proj}}\) is not a reliable short-window estimate of
the true MRW intermittency parameter.

The practical implication is conservative. CMIN-SR can organize evidence about
scaling, curvature, projection residuals, and instability. It should not be
used to announce an MRW mechanism merely because a finite spectrum can be fit
by the MRW family.
